\section{Introduction}
\label{sec:intro}

Streaming services recommend movies one at a time, and every suggestion changes what happens next: a user who rejects several suggestions in a row may simply leave. Most recommenders, however, are trained as static predictors that score each movie once per user. Our previous workshops followed that approach on the MovieLens 32M dataset~\cite{harper2015movielens}, which holds 32{,}000{,}204 ratings given by 200{,}948 users to 87{,}585 movies. Workshop~1 trained FunkSVD matrix factorization with user and item biases~\cite{koren2009matrix}, which reached a test RMSE of 0.8041 with 50 latent factors. Workshop~2 trained a hybrid neural collaborative filtering model (HybridNeuMF) that lowered the test RMSE to 0.7994 and reached an NDCG@10 of 0.7663 when each test user's liked movies were ranked among sampled unrated ones. Both models produce one fixed ranking per user and ignore the feedback the user gives while browsing.

Reinforcement learning (RL) formulates the same task as sequential decision making~\cite{sutton2018reinforcement,sierra2026rl}. The recommender becomes an agent that observes the state of a session, recommends an item, receives a reward from the user's response, and learns a policy that maximizes the cumulative reward of the session instead of the accuracy of each prediction. Such an agent can, in principle, adapt to feedback inside a session and avoid streaks of poor recommendations that end it. It also faces the exploration--exploitation trade-off: it must try uncertain actions to learn their value while exploiting what it already knows~\cite{sutton2018reinforcement}. Among value-based methods, Q-learning~\cite{watkins1992q} learns action values with temporal-difference backups, the Deep Q-Network (DQN)~\cite{mnih2015human} extends it to high-dimensional states with a replay buffer and a target network, Double DQN~\cite{van2016deep} reduces the overestimation caused by maximizing over noisy estimates, and the dueling architecture~\cite{wang2016dueling} separates the value of a state from the advantage of each action. Afsar et al.~\cite{afsar2022reinforcement} survey RL-based recommenders. Zhao et al.~\cite{zhao2018recommendations} trained a deep Q-learning recommender that learns from both positive and negative feedback, and SlateQ~\cite{ie2019slateq} decomposes the value of a recommendation set into item-wise values so that Q-learning stays tractable when many items compete.

Applying RL to logged ratings raises three challenges. First, the agent cannot interact with real users, so it must learn from logged data or from a simulator built on it; offline RL suffers from the distribution shift between the logged behavior and the learned policy~\cite{levine2020offline}. Second, a simulator can leak information that would not be available at decision time. Kaufman et al.~\cite{kaufman2012leakage} describe leakage as information about the target that enters the training data although it would not be legitimately available, and show that it yields models that look excellent offline and fail in use. Third, once sessions are restricted to known users and movies after a temporal cut, they are few and short, which limits what a value function can learn. The first version of this workshop failed on the second point: its simulator exposed the hidden labels in the observation, so its results were artifacts (Section~\ref{sec:challenges}).

This paper reports the rebuilt study. We cast sequential top-$K$ recommendation as a finite-horizon Markov decision process (MDP) whose sessions are the real ratings that known users gave after the training period, implement it as a Gymnasium environment~\cite{towers2024gymnasium} whose observation is computed only from training-split models and observed feedback, and train a Dueling Double DQN (D3QN) agent with two Q-network architectures. The agent competes with the Workshop~1 and Workshop~2 models acting greedily in the same environment and with a copy of itself trained with $\gamma = 0$, which isolates the value of planning; every comparison uses the same held-out sessions and paired statistical tests. The agent matches the supervised rankers in return but not in ranking quality, which shows how little room a log-replay simulator leaves for sequential decisions.
Section~\ref{sec:methods} describes the data, the MDP, the agent and the evaluation protocol, Section~\ref{sec:results} reports and discusses the results, and Section~\ref{sec:conclusions} concludes. The code, data artifacts and checkpoints are available at \url{https://github.com/AndresRuiz14502/Machine_Learning}.
